\documentclass[11pt,a4paper]{article}
\usepackage[margin=2.45cm]{geometry}
\usepackage[T1]{fontenc}
\usepackage{lmodern,microtype}
\usepackage{amsmath,amssymb,amsfonts,booktabs,graphicx,float,hyperref,xcolor,enumitem,array}
\usepackage{algorithm,algpseudocode}
\hypersetup{colorlinks=true,linkcolor=blue!50!black,citecolor=blue!50!black,urlcolor=blue!50!black}
\newcommand{\norm}[1]{\left\lVert#1\right\rVert_2}
\newcommand{\mean}{\boldsymbol{\mu}}
\newcommand{\upd}{\boldsymbol{\Delta}}
\newcommand{\mal}{\boldsymbol{m}}
\newcommand{\dir}{\boldsymbol{d}}
\newcommand{\agg}{\mathcal{A}}
\title{\textbf{Subspace-Constrained Byzantine Attacks}\\[5pt]\large Weighted PCA Directions under Robust Aggregation}
\author{Po-Kai Huang\\TUM PREP 2026\\[4pt]\small Supervisors: Sena Ergisi and Rawad Bitar}
\date{October 2026}
\begin{document}
\maketitle
\begin{abstract}
Byzantine model-poisoning attacks challenge the robustness of federated learning by allowing compromised clients to submit arbitrary model updates. Existing attacks such as Min-Max and Min-Sum constrain the magnitude of malicious updates but typically operate along predefined perturbation directions. This report investigates whether a low-dimensional, data-dependent search space constructed from benign model-update geometry offers effective alternative directions. We develop a weighted principal-component approach that recomputes the benign subspace each communication round, combines its leading components using standard-deviation weights, enumerates component-wise orientations, and selects a candidate using an aggregation-shift proxy under a Min-Max distance constraint. We evaluate two- and three-component variants on CIFAR-10 with ten clients, three Byzantine participants, five aggregation rules, and optional nearest-neighbor mixing (NNM). In the reported NNM comparisons, the better PCA variant achieves lower mean final accuracy than the best non-sign-flip classical baseline for four of five aggregation rules. The advantage is not universal without NNM. We discuss limitations arising from incomplete runs, baseline construction, and discrepancies between the candidate-selection proxy and the full defended pipeline.
\end{abstract}

\section{Introduction}
Federated learning (FL) coordinates training across clients that communicate parameter updates to a central server. The server aggregates these updates to obtain a new global model. Under a Byzantine threat model, some clients may replace their updates with arbitrary vectors, potentially corrupting learning even when the server employs robust aggregation.

A useful representation of a poisoning update separates its direction from its magnitude. Classical distance-constrained attacks often choose a predefined direction and search for the largest admissible displacement. This motivates the question: \emph{Can directions derived from the geometry of benign updates offer useful alternatives under robust aggregation?}

The present work explores a deliberately restricted answer. Principal component analysis (PCA) does not guarantee a harmful direction; rather, it supplies a compact, data-dependent basis in which a finite collection of weighted directions can be evaluated. Our contributions are (i) a weighted PCA and component-wise sign-enumeration attack, (ii) an empirical comparison across five aggregation rules with and without NNM, and (iii) an analysis of the sensitivity of attack outcomes to defense preprocessing and experimental evaluation choices.

\section{Background and Related Work}
\subsection{Federated model updates and Byzantine behavior}
At round $t$, let $\boldsymbol{\theta}^{(t)}\in\mathbb{R}^D$ denote global model parameters. Benign client $i$ trains locally and returns a model update
\begin{equation}
\upd_i^{(t)}=\boldsymbol{\theta}_i^{(t,\mathrm{local})}-\boldsymbol{\theta}^{(t)}.
\end{equation}
The server forms $\upd_G^{(t)}=\agg(\upd_1^{(t)},\ldots,\upd_n^{(t)})$ and updates $\boldsymbol{\theta}^{(t+1)}=\boldsymbol{\theta}^{(t)}+\upd_G^{(t)}$. Unlike ordinary stochastic gradients, the vectors analyzed here are \emph{model parameter differences} following local training.

The tested aggregation methods are FedAvg, coordinate-wise median, Krum, Multi-Krum, and a median-centered coordinate-selection variant labeled TrimmedMean in the implementation. FedAvg takes the unweighted mean of submitted client updates. Krum chooses a single update using distances to its $n-f-2$ nearest other updates; Multi-Krum averages the $n-f-2$ lowest-scoring updates. The implemented TrimmedMean variant averages the $n-2f$ coordinate values closest to the coordinate-wise median, rather than discarding the $f$ largest and $f$ smallest values. Its results should therefore not be presented as results for standard trimmed mean.

\subsection{Reference attacks}
ALIE exploits coordinate-wise benign statistics to construct a malicious update \cite{baruch2019}. Min-Max and Min-Sum constrain the attack relative to benign pairwise distances \cite{shejwalkar2021}. The implementation used here also includes a sign-flip baseline, which negates the benign mean. The classical Min-Max and Min-Sum implementations use a direction based on the negative sign of the benign mean; the PCA attack replaces that direction with weighted principal-component combinations.

\subsection{Nearest-neighbor mixing}
NNM is a preprocessing step applied before aggregation. For each submitted vector, it selects the $n-f$ nearest vectors \emph{including itself} and replaces the original vector with their average. Thus the effective aggregation pipeline can be written $\agg(\mathcal{N}(U))$, where $\mathcal{N}$ denotes NNM and $U$ is the collection of submitted updates. This transformation can change both the geometry and relative influence of malicious submissions.

\section{Threat Model and Attack Objective}
Let $B=\{\boldsymbol{b}_1,\ldots,\boldsymbol{b}_{n-f}\}$ be the current round's benign updates, and let $f$ Byzantine clients collude by submitting an identical vector $\mal$. We assume a strong, information-rich attacker that observes \emph{all current benign updates} and knows the target aggregation rule. This is not a limited-information or black-box attack model.

Define the benign mean and maximum benign pairwise distance as
\begin{equation}
\mean=\frac{1}{n-f}\sum_{i=1}^{n-f}\boldsymbol{b}_i,\qquad
R_B=\max_{i,j}\norm{\boldsymbol{b}_i-\boldsymbol{b}_j}.
\end{equation}
For a unit candidate direction $\dir$, the attack constructs
\begin{equation}\label{eq:mal}
\mal(\dir)=\mean-\gamma(\dir)\dir,
\end{equation}
where $\gamma$ is chosen as large as possible subject to the implemented Min-Max feasibility condition
\begin{equation}\label{eq:constraint}
\max_i\norm{\mal(\dir)-\boldsymbol{b}_i}\leq R_B.
\end{equation}
This restricts malicious updates to a benign-distance-defined feasible region, but does not guarantee that a robust aggregator will accept or favor them.

\section{Weighted PCA Direction Search}
\subsection{Benign-update subspace}
Write $q=n-f$ and form the centered matrix $X\in\mathbb{R}^{q\times D}$ with row $i$ equal to $(\boldsymbol{b}_i-\mean)^\top$. Compute a thin singular value decomposition
\begin{equation}
X=U\Sigma V^\top,\qquad \lambda_j=\frac{\sigma_j^2}{q-1}.
\end{equation}
The columns $\boldsymbol{v}_j$ of $V$ are principal directions, with $K\leq q-1$. We use the first $K\in\{2,3\}$ directions. PCA is computed on \emph{client model updates}, not on the input images, and is recomputed each round.

\subsection{Weighted directions and orientation search}
Set $w_j=\sqrt{\lambda_j}$, corresponding to the standard deviation along each principal component. For a sign vector $\boldsymbol{s}\in\{-1,+1\}^K$, construct
\begin{equation}\label{eq:direction}
\dir(K,\boldsymbol{s})=\frac{\sum_{j=1}^{K}s_jw_j\boldsymbol{v}_j}{\norm{\sum_{j=1}^{K}s_jw_j\boldsymbol{v}_j}}.
\end{equation}
There are $2^K$ sign configurations: four for $K=2$ and eight for $K=3$. These are \emph{discrete weighted combinations}, not a continuous optimization over all vectors in the PCA subspace. Because all sign vectors are enumerated, using $\mean+\gamma\dir$ instead of Eq.~\eqref{eq:mal} would produce the same candidate direction set for this experiment.

\subsection{Magnitude search and candidate scoring}
For each sign configuration, binary search determines a feasible displacement magnitude $\gamma$ satisfying Eq.~\eqref{eq:constraint}. The candidate is then scored using the implemented aggregation-space displacement
\begin{equation}\label{eq:score}
J(\boldsymbol{s})=\norm{\agg_f(B\cup M(\boldsymbol{s}))-\agg_f(B)},
\qquad M(\boldsymbol{s})=\{\mal(\dir(K,\boldsymbol{s}))\}^{f},
\end{equation}
where $\agg_f$ indicates the configured aggregator with its malicious-count parameter set to $f$. The second term is the aggregator output on benign-only updates under that same configuration; it is \emph{not necessarily the benign arithmetic mean}. The selected sign vector is
\begin{equation}
\boldsymbol{s}^{\star}=\arg\max_{\boldsymbol{s}\in\{-1,+1\}^{K}}J(\boldsymbol{s}).
\end{equation}
The objective is a proxy for disruption in update space, not test accuracy. Importantly, candidate scoring \emph{does not include NNM}, even when the eventual server pipeline uses it. Consequently, the method is aggregation-aware but not fully NNM-aware.

\begin{algorithm}[t]
\caption{Weighted PCA attack at one communication round}
\begin{algorithmic}[1]
\Require Benign updates $B$, aggregator $\agg_f$, Byzantine count $f$, number of PCs $K$
\State Compute benign mean $\mean$ and centered matrix $X$
\State Obtain top $K$ PCA directions $\boldsymbol{v}_j$ and weights $w_j=\sqrt{\lambda_j}$
\State Compute maximum benign pairwise distance $R_B$
\For{each $\boldsymbol{s}\in\{-1,+1\}^K$}
  \State Construct $\dir(K,\boldsymbol{s})$ using Eq.~\eqref{eq:direction}
  \State Binary-search the largest feasible $\gamma$ under Eq.~\eqref{eq:constraint}
  \State Set $\mal(\boldsymbol{s})=\mean-\gamma\dir(K,\boldsymbol{s})$
  \State Evaluate $J(\boldsymbol{s})$ using Eq.~\eqref{eq:score}
\EndFor
\State \Return $f$ copies of $\mal(\boldsymbol{s}^{\star})$
\end{algorithmic}
\end{algorithm}

\section{Experimental Setup}
\subsection{Dataset, model, and training}
Experiments use CIFAR-10 with a simplified convolutional neural network comprising three convolution--ReLU--max-pooling blocks (32, 64, and 128 channels), a 128-unit hidden layer, and ten output classes. Data are distributed non-IID across clients using a Dirichlet partition with concentration parameter $\alpha=0.5$. There are ten clients per round, of which three are Byzantine; all ten participate. Training runs for 50 rounds with one local epoch per round, batch size 256, SGD momentum 0.9, and learning rate 0.05 for rounds 1--33 followed by 0.005 for rounds 34--50. The three main random seeds are 40, 60, and 11.

\subsection{Comparison protocol}
We compare weighted PCA with $K=2$ and $K=3$ against ALIEOptim, Min-Max, Min-Sum, and Sign Flip across five aggregators, each both directly (raw) and after NNM. The primary metric is CIFAR-10 test accuracy, with smaller values indicating greater degradation. Each seed's reported summary averages rounds 46--50; tables show the mean and sample standard deviation across three seeds. The reported no-attack baseline trains only the seven clients designated benign in the attack experiments, rather than ten benign clients. Under the no-attack NNM setting, $f=0$, so NNM averages all seven updates into the same vector before aggregation.

\subsection{Incomplete runs and imputation}
The archived experiment inventory contains 210 main CSVs; 176 have all 51 records (including round zero), and 34 ended early. The analysis notebook fills missing late-round accuracy values with 0.10 before computing the final-five-round statistics. Accordingly, a reported 10\% value may be either measured or imputed and cannot automatically be interpreted as convergence to chance accuracy. We reproduce the presentation's tabulated summaries, but conclusions should be interpreted with this handling in mind. The presentation attributes Sign Flip failures to logit overflow; the available error summary does not establish a cause for every interrupted run.

\section{Results and Analysis}
\subsection{Final accuracy without NNM}
\begin{table}[htbp]\centering\caption{Mean CIFAR-10 accuracy (\%) over the final five rounds, \textbf{without} NNM; mean $\pm$ sample standard deviation over three seeds. Values reproduce the presentation and include its missing-round imputation policy.}\label{tab:raw}\scriptsize\setlength{\tabcolsep}{3pt}
\begin{tabular}{lrrrrrrr}\toprule
Aggregator & No attack & ALIE & Min-Max & Min-Sum & Sign Flip & PCA $K=2$ & PCA $K=3$\\\midrule
FedAvg & 64.31$\pm$1.26 & 40.33$\pm$10.86 & 21.33$\pm$19.63 & 36.82$\pm$23.28 & 49.90$\pm$0.94 & 49.17$\pm$3.20 & 49.55$\pm$3.81\\
Median & 55.94$\pm$3.11 & 21.29$\pm$7.64 & 39.02$\pm$2.78 & 41.52$\pm$3.41 & 23.56$\pm$1.25 & 33.82$\pm$6.65 & 34.50$\pm$2.86\\
TrimmedMean$^\dagger$ & 64.34$\pm$1.29 & 15.37$\pm$4.94 & 38.72$\pm$3.13 & 40.37$\pm$3.99 & 21.38$\pm$3.97 & 32.29$\pm$2.19 & 27.44$\pm$7.94\\
Krum & 25.26$\pm$17.69 & 10.00$\pm$0.00 & 10.00$\pm$0.00 & 9.99$\pm$0.01 & 19.06$\pm$15.69 & 10.00$\pm$0.00 & 10.00$\pm$0.00\\
Multi-Krum & 61.01$\pm$1.60 & 13.79$\pm$6.57 & 18.29$\pm$14.37 & 10.00$\pm$0.00 & 10.00$\pm$0.00 & 10.00$\pm$0.00 & 10.00$\pm$0.00\\\bottomrule
\end{tabular}
\par\smallskip\raggedright\footnotesize $^\dagger$Median-centered coordinate-selection variant, not standard trimmed mean.
\end{table}

Without NNM, PCA does not dominate classical attacks. For example, ALIE produces 15.37\% on the implemented TrimmedMean variant, whereas PCA with $K=2$ and $K=3$ yields 32.29\% and 27.44\%. The no-attack Krum baseline is itself unstable (25.26 $\pm$ 17.69\%), complicating attribution of poor attacked accuracy to the attack alone.

\subsection{Final accuracy with NNM}
\begin{table}[htbp]\centering\caption{Mean CIFAR-10 accuracy (\%) over the final five rounds, \textbf{with} NNM; mean $\pm$ sample standard deviation over three seeds. Imputation caveats apply.}\label{tab:nnm}\scriptsize\setlength{\tabcolsep}{3pt}
\begin{tabular}{lrrrrrrr}\toprule
Aggregator & No attack & ALIE & Min-Max & Min-Sum & Sign Flip & PCA $K=2$ & PCA $K=3$\\\midrule
FedAvg & 64.28$\pm$1.22 & 48.94$\pm$6.15 & 36.39$\pm$22.99 & 34.19$\pm$21.09 & 40.14$\pm$7.15 & 47.94$\pm$2.94 & 46.36$\pm$5.10\\
Median & 64.32$\pm$1.32 & 37.29$\pm$12.17 & 39.91$\pm$25.91 & 34.78$\pm$21.48 & 10.00$\pm$0.00 & 29.14$\pm$13.82 & 31.04$\pm$2.27\\
TrimmedMean$^\dagger$ & 64.33$\pm$1.30 & 38.50$\pm$11.30 & 40.54$\pm$26.46 & 24.68$\pm$16.44 & 10.00$\pm$0.00 & 22.58$\pm$9.04 & 21.92$\pm$6.95\\
Krum & 64.31$\pm$1.27 & 35.56$\pm$11.29 & 39.80$\pm$25.82 & 23.26$\pm$22.98 & 11.58$\pm$2.74 & 18.44$\pm$6.74 & 16.96$\pm$6.12\\
Multi-Krum & 64.29$\pm$1.25 & 41.21$\pm$11.42 & 39.89$\pm$25.89 & 20.78$\pm$18.67 & 10.00$\pm$0.00 & 27.16$\pm$6.43 & 18.34$\pm$7.28\\\bottomrule
\end{tabular}
\par\smallskip\raggedright\footnotesize $^\dagger$Median-centered coordinate-selection variant, not standard trimmed mean.
\end{table}

Under NNM, the better of the two PCA variants obtains lower mean accuracy than the strongest non-sign-flip classical baseline for Median (29.14\% versus 34.78\%), the implemented TrimmedMean variant (21.92\% versus 24.68\%), Krum (16.96\% versus 23.26\%), and Multi-Krum (18.34\% versus 20.78\%). FedAvg is an exception: Min-Sum yields 34.19\%, compared with the best PCA result of 46.36\%. Thus the ``four of five'' comparison is a per-aggregator \emph{best-of-two-$K$} observation, not a claim that each PCA variant wins in four settings, nor a statistical significance result. Sign Flip is shown for transparency but should be analyzed separately because its numerical failures complicate comparison with stable degraded learning.

\subsection{Update-space diagnostics}
The presentation includes per-round attack aggregation shift, malicious-to-benign-mean Euclidean distance, and cosine similarity for the implemented TrimmedMean variant, with and without NNM. The displayed NNM traces suggest that PCA can produce relatively large aggregate shifts without consistently having larger malicious-to-benign distances than classical attacks. PCA malicious vectors also exhibit greater directional alignment with the benign mean than several reference attacks in the illustrated rounds. These observations support the plausibility of geometry-aligned disruption but do not establish a causal explanation. The diagnostics compare the attacked aggregate with a separate clean aggregation pipeline, whereas the selection score in Eq.~\eqref{eq:score} is computed without NNM and with a differently configured benign reference.

\section{Limitations and Future Work}
Several limitations constrain interpretation. First, experiments cover only CIFAR-10, one non-IID partition concentration, one Byzantine fraction, two values of $K$, and three seeds. Second, a significant subset of runs ended early; imputation with 10\% accuracy may materially affect means and rankings. Third, the benign-only baseline uses seven active clients rather than a ten-benign-client reference, and the reported TrimmedMean implementation differs from the conventional definition. Fourth, maximizing an aggregation-shift proxy does not guarantee a corresponding decrease in test accuracy, and the proxy does not model NNM when NNM is used in deployment. Fifth, the strong attacker observes all benign updates, which may not be realistic in less permissive settings.

Future work should first separate completed runs from numerical failures and report sensitivity to missing-data handling, then validate the conventional trimmed mean and a consistent clean baseline. Methodologically, the next step is to optimize direction coefficients within the PCA subspace rather than enumerate only signs, and to score candidates against the \emph{entire} preprocessing-plus-aggregation pipeline. Additional datasets, client partitions, Byzantine fractions, and repeated seeds would help establish generality.

\section{Conclusion}
Weighted PCA provides a structured, low-dimensional source of Byzantine attack directions derived from current benign model updates. In the reported CIFAR-10 experiments, a discrete sign search over two or three principal components yields effective attacks under several robust aggregation rules. Its relative advantage is conditional: with NNM, the better PCA variant outperforms the best non-sign-flip classical baseline in four of five mean-accuracy comparisons, while without NNM the results are mixed. These findings motivate more adaptive, defense-aware direction optimization, while also emphasizing the importance of accurate implementation descriptions and careful handling of interrupted experiments.

\begin{thebibliography}{9}
\bibitem{baruch2019} G. Baruch, M. Baruch, and Y. Goldberg. ``A Little Is Enough: Circumventing Defenses for Distributed Learning.'' \emph{Advances in Neural Information Processing Systems}, 2019.
\bibitem{shejwalkar2021} V. Shejwalkar and A. Houmansadr. ``Manipulating the Byzantine: Optimizing Model Poisoning Attacks and Defenses for Federated Learning.'' \emph{Network and Distributed System Security Symposium (NDSS)}, 2021.
\end{thebibliography}
\end{document}
